\begin{thebibliography}{99}
\small
\setlength{\itemsep}{0.45em}

\bibitem{raghavan2020} Raghavan, M., Barocas, S., Kleinberg, J. y Levy, K. (2020). \textit{Mitigating Bias in Algorithmic Hiring: Evaluating Claims and Practices}. Proceedings of FAT* 2020. \url{https://doi.org/10.1145/3351095.3372828}. Versión abierta: \url{https://arxiv.org/abs/1906.09208}.

\bibitem{manning2008} Manning, C. D., Raghavan, P. y Schütze, H. (2008). \textit{Introduction to Information Retrieval}. Cambridge University Press. Sección sobre ponderación TF--IDF. \url{https://nlp.stanford.edu/IR-book/html/htmledition/tf-idf-weighting-1.html}.

\bibitem{reimers2019} Reimers, N. y Gurevych, I. (2019). Sentence-BERT: Sentence Embeddings using Siamese BERT-Networks. \textit{Proceedings of EMNLP-IJCNLP}, 3982--3992. \url{https://doi.org/10.18653/v1/D19-1410}.

\bibitem{suresh2021} Suresh, H. y Guttag, J. V. (2021). \textit{A Framework for Understanding Sources of Harm throughout the Machine Learning Life Cycle}. Versión revisada de arXiv:1901.10002, 1 de diciembre de 2021; trabajo asociado a EAAMO 2021. \url{https://arxiv.org/abs/1901.10002}.

\bibitem{dearteaga2019} De-Arteaga, M., Romanov, A., Wallach, H., Chayes, J., Borgs, C., Chouldechova, A., Geyik, S., Kenthapadi, K. y Kalai, A. T. (2019). \textit{Bias in Bios: A Case Study of Semantic Representation Bias in a High-Stakes Setting}. Proceedings of FAT* 2019. \url{https://doi.org/10.1145/3287560.3287572}. Versión abierta: \url{https://arxiv.org/abs/1901.09451}.

\bibitem{hardt2016} Hardt, M., Price, E. y Srebro, N. (2016). \textit{Equality of Opportunity in Supervised Learning}. Advances in Neural Information Processing Systems, 29. \url{https://arxiv.org/abs/1610.02413}.

\bibitem{kleinberg2017} Kleinberg, J., Mullainathan, S. y Raghavan, M. (2017). \textit{Inherent Trade-Offs in the Fair Determination of Risk Scores}. Proceedings of ITCS 2017. Preprint de 2016: \url{https://arxiv.org/abs/1609.05807}.

\bibitem{kusner2017} Kusner, M. J., Loftus, J. R., Russell, C. y Silva, R. (2017). \textit{Counterfactual Fairness}. Advances in Neural Information Processing Systems, 30. \url{https://arxiv.org/abs/1703.06856}.

\bibitem{saldivar2025} Saldivar, J., Gatzioura, A. y Castillo, C. (2025). \textit{Synthetic CVs To Build and Test Fairness-Aware Hiring Tools}. Preprint, arXiv:2508.21179. \url{https://arxiv.org/abs/2508.21179}.

\bibitem{kamiran2012} Kamiran, F. y Calders, T. (2012). Data preprocessing techniques for classification without discrimination. \textit{Knowledge and Information Systems}, 33, 1--33. \url{https://doi.org/10.1007/s10115-011-0463-8}.

\bibitem{agarwal2018} Agarwal, A., Beygelzimer, A., Dudík, M., Langford, J. y Wallach, H. (2018). A Reductions Approach to Fair Classification. \textit{Proceedings of Machine Learning Research}, 80, 60--69. \url{https://proceedings.mlr.press/v80/agarwal18a.html}.

\bibitem{gebru2021} Gebru, T., Morgenstern, J., Vecchione, B., Vaughan, J. W., Wallach, H., Daumé III, H. y Crawford, K. (2021). Datasheets for Datasets. \textit{Communications of the ACM}. \url{https://doi.org/10.1145/3458723}. Preprint: \url{https://arxiv.org/abs/1803.09010}.

\bibitem{mitchell2019} Mitchell, M., Wu, S., Zaldivar, A., Barnes, P., Vasserman, L., Hutchinson, B., Spitzer, E., Raji, I. D. y Gebru, T. (2019). \textit{Model Cards for Model Reporting}. Proceedings of FAT* 2019. \url{https://doi.org/10.1145/3287560.3287596}. Versión abierta: \url{https://arxiv.org/abs/1810.03993}.

\bibitem{aiact2024} Parlamento Europeo y Consejo de la Unión Europea (2024). \textit{Reglamento (UE) 2024/1689, de 13 de junio de 2024, por el que se establecen normas armonizadas en materia de inteligencia artificial}. Artículo 6, capítulo III y anexo III. Texto original: \url{https://eur-lex.europa.eu/eli/reg/2024/1689/oj?locale=es}.

\bibitem{comision2026} Comisión Europea (2026). \textit{AI Act}. Shaping Europe's digital future. Información oficial sobre el marco regulatorio y su calendario de aplicación. Consultado el 13 de octubre de 2026. \url{https://digital-strategy.ec.europa.eu/en/policies/regulatory-framework-ai}.

\bibitem{rgpd2016} Parlamento Europeo y Consejo de la Unión Europea (2016). \textit{Reglamento (UE) 2016/679, de 27 de abril de 2016 (Reglamento general de protección de datos)}. Artículos 5, 6, 9 y 22. \url{https://www.boe.es/buscar/doc.php?id=DOUE-L-2016-80807}.

\bibitem{ley152022} Jefatura del Estado (2022). \textit{Ley 15/2022, de 12 de julio, integral para la igualdad de trato y la no discriminación}. Artículos 9 y 23. Texto consolidado consultado el 13 de octubre de 2026. \url{https://www.boe.es/buscar/act.php?id=BOE-A-2022-11589}.

\bibitem{wilson2024} Wilson, K. y Caliskan, A. (2024). \textit{Gender, Race, and Intersectional Bias in Resume Screening via Language Model Retrieval}. Proceedings of AIES 2024. \url{https://doi.org/10.1609/aies.v7i1.31748}. Se ha consultado también la versión de arXiv revisada en agosto de 2026, que incorpora una descripción de erratas: \url{https://arxiv.org/abs/2407.20371v3}.

\bibitem{sariola2026} Sariola, D., Button, P., Culotta, A. y Mattei, N. (2026). The Illusion of Fairness: Auditing Fairness Interventions in Algorithmic Hiring with Audit Studies. \textit{Proceedings of the AAAI Conference on Artificial Intelligence}, 40(46), 39191--39200. \url{https://doi.org/10.1609/aaai.v40i46.41267}.

\bibitem{sklearnapi} Scikit-learn developers (2026). \textit{API reference: TfidfVectorizer, LogisticRegression y LinearSVC}. Documentación oficial, consultada el 13 de octubre de 2026. \url{https://scikit-learn.org/stable/modules/generated/sklearn.feature_extraction.text.TfidfVectorizer.html}; \url{https://scikit-learn.org/stable/modules/generated/sklearn.linear_model.LogisticRegression.html}; \url{https://scikit-learn.org/stable/modules/generated/sklearn.svm.LinearSVC.html}.

\bibitem{repositorio2026} Rajo Vázquez, J. (2026). \textit{Evaluación y mitigación de sesgos en reclutamiento: corpus sintético, protocolo y resultados reproducibles} (versión 2.0.0) [Software y datos]. GitHub. \url{https://github.com/jordanrajo-svg/tfm-sesgos-reclutamiento}. Consulta: 13 de octubre de 2026. Incluye configuración, dependencias, generador, entrenamiento, predicciones, verificador y figuras. Reproducción desde la raíz del repositorio: \texttt{python -m pip install -r requirements.txt}, seguido de \texttt{python reproducir.py}. Solo verificación y gráficos: \texttt{python reproducir.py --solo-analisis}. La identificación de la versión se completa con el historial de commits y las huellas del registro de ejecución.

\bibitem{omnibus2026} Comisión Europea (2026). \textit{AI Omnibus enters into force}. 27 de julio de 2026. \url{https://digital-strategy.ec.europa.eu/en/news/ai-omnibus-enters-force}. Consulta: 13 de octubre de 2026.
\end{thebibliography}